
% --- SLIDE 2: SMBO ---
\begin{frame}{Sequential Model-Based Optimisation: The Idea}

    \centering
    \begin{tikzpicture}[
        b/.style={rectangle, draw, rounded corners=3pt, minimum height=1.0cm, text width=2.5cm,
                  align=center, font=\scriptsize},
        arr/.style={-{Latex[length=2.2mm]}, thick, gray!60!black}
    ]
        \node[b, fill=blue!12]   (hist) at (0,0)     {\textbf{History}\\$(x_i, y_i)$ so far};
        \node[b, fill=green!12]  (surr) at (3.4,0)   {\textbf{Surrogate model}\\cheap guess of\\$y$ given $x$};
        \node[b, fill=orange!12] (acq)  at (6.8,0)   {\textbf{Acquisition}\\where is the next\\trial worth most?};
        \node[b, fill=red!10]    (eval) at (10.2,0)  {\textbf{Evaluate}\\actually train\\the model};

        \draw[arr] (hist) -- (surr);
        \draw[arr] (surr) -- (acq);
        \draw[arr] (acq)  -- (eval);
        \draw[arr] (eval.south) -- ++(0,-0.85) -| node[pos=0.25, below, font=\scriptsize] {append} (hist.south);
    \end{tikzpicture}

    \vspace{1.0em}
    \begin{itemize}
        \item Training one configuration is \emph{expensive}; evaluating a surrogate is \emph{free}. So spend
              compute on the surrogate to decide where to spend compute on training.
        \item The acquisition function balances \textbf{exploitation} (near the best point seen) against
              \textbf{exploration} (where the surrogate is uncertain).
        \item Gaussian-process Bayesian optimisation is the textbook version. \textbf{TPE} is the version
              that scales to messy, conditional, mixed-type search spaces --- which is what real projects have.
    \end{itemize}

\end{frame}
